\chapter{Пробный полный цикл реконструкции}

Общий протокол должен быть проверен на одном заранее ограниченном классе
данных. Цель пробного цикла --- проверить работоспособность процедуры, а не
объявить выбранный пример фундаментальным описанием мира.

\section{Выбор класса данных}

Рассмотрим условный набор
\begin{equation}
 \mathcal D^{(1)}=
 \{\rho(\Lambda),\Delta_\alpha,m_a,C_{a\alpha}\},
\end{equation}
где $\rho$ обозначает спектральную плотность, $\Delta_\alpha$ --- различие
глобальных фазовых ветвей, $m_a$ --- внутренние уровни, а $C_{a\alpha}$ ---
смешанные коэффициенты отклика.

Этот набор выбран потому, что содержит каркасный, фазовый, внутренний и
смешанный типы данных. Его символы не отождествляются автоматически с
конкретными наблюдаемыми Стандартной модели.

\section{Выделение устойчивого содержания}

В $\mathcal D^{(1)}_{\mathrm{stab}}$ включаются только закономерности,
сохраняющиеся при изменении сетки, базиса и допустимой нормировки. Отдельное
число без устойчивой структуры исключается из ядра пробного цикла.

\section{Секторная атрибуция}

Предварительная карта имеет вид
\begin{align}
 \rho(\Lambda)&\longmapsto\mathrm{car},\\
 \Delta_\alpha&\longmapsto\mathrm{phase},\\
 m_a&\longmapsto\mathrm{int},\\
 C_{a\alpha}&\longmapsto\mathrm{mix}.
\end{align}

Эта карта является гипотезой до проверки деформационной чувствительности.
Если, например, $C_{a\alpha}$ раскладывается на независимые фазовую и
внутреннюю части без перекрёстного остатка, смешанный статус отзывается.

\section{Обратная реконструкция}

Из карты ищется минимальная структура, способная породить все четыре типа
одним правилом. Успех требует общего носителя и совместимого семейства
операторов. Четыре независимых формулы не считаются реконструкцией единого
источника.

\section{Прямое следствие}

После обратной реконструкции модель должна ограничить хотя бы один новый
коэффициент, знак, ранг или допустимую форму смешения. Если новое ограничение
отсутствует, цикл остаётся способом классификации уже известных данных.

\chapter{Пределы применимости и критерии опровержения}

Программа обязана заранее перечислять обстоятельства, при которых её основные
утверждения теряют силу.

\section{Опровержение единого источника}

Гипотеза $\mathfrak S$ не проходит проверку в объявленном классе, если:
\begin{itemize}
 \item разные языки требуют несовместимых инвариантных ядер;
 \item общий носитель существует только после независимой настройки секторов;
 \item переходы между редукциями неоднозначны без дополнительных данных;
 \item новые наблюдаемые не получают ограничений.
\end{itemize}

\section{Опровержение секторной карты}

Карта режимов отзывается, если классификация существенно зависит от выбора
координат, нормировки или порядка обработки данных. Граница областей должна
быть следствием модели, а не рисунком, проведённым после вычисления.

\section{Опровержение модели без опровержения программы}

Провал конкретного носителя, оператора или функционала не опровергает всю
методологию S2T. Он закрывает только явно объявленную архитектуру. Более общий
вывод допустим лишь после классификации всего заявленного класса кандидатов.

\section{Уровни отрицательного результата}

Следует различать:
\begin{enumerate}
 \item численный провал выбранного набора параметров;
 \item структурный запрет для семейства моделей;
 \item классификационный запрет для объявленного класса;
 \item универсальную теорему невозможности.
\end{enumerate}

Каждый следующий уровень требует отдельного доказательства.

\chapter{Иерархия конкурирующих моделей}

Один успешный сравнительный тест не задаёт устойчивого порядка моделей.
Необходимо проверить их на нескольких различающихся классах наблюдаемых.

\section{Матрица оценивания}

Для модели $M$ и класса $\mathcal C_k$ вводится ведомость
\begin{equation}
 \mathcal Q_k(M)=
 \bigl(q_{\mathrm{reg}},q_{\mathrm{econ}},q_{\mathrm{stab}},
 q_{\mathrm{trans}},q_{\mathrm{pred}}\bigr),
\end{equation}
где учитываются режимная совместимость, экономия, устойчивость,
переносимость и различающая сила.

Компоненты не обязаны сворачиваться в одно число. Скалярная сумма допустима
только после независимого вывода весов.

\section{Локальное и глобальное предпочтение}

Модель может быть лучшей на одном классе и худшей на другом. Поэтому
различаются локальное предпочтение $M\succ_k M'$ и глобальное предпочтение,
устойчивое на семействе классов.

\section{Статусы моделей}

После многоклассовой проверки модель получает один из статусов:
\begin{itemize}
 \item фундаментальный кандидат;
 \item эффективное приближение;
 \item переходная конструкция;
 \item вычислительный инструмент;
 \item закрытая архитектура.
\end{itemize}

Проигрыш в фундаментальном отборе не уничтожает локальную вычислительную
полезность, но запрещает расширять область утверждений.

\section{Защита от смены правил}

Классы данных, показатели качества и порядок разрешения ничьих фиксируются до
сравнения. Добавление нового показателя после проигрыша модели открывает новое
исследование, а не исправляет прежний результат.

\chapter{Граница между программой и строгой теорией}

Исследовательская программа становится теорией не благодаря объёму текста, а
после замыкания цепи от определения до независимой проверки.

\section{Необходимые уровни}

Минимальная цепь включает:
\begin{enumerate}
 \item точное определение носителя и пространства состояний;
 \item оператор или действие без скрытых секторных параметров;
 \item доказательство существования требуемого режима;
 \item проверку устойчивости и области применимости;
 \item вывод наблюдаемого следствия;
 \item независимый воспроизводимый контроль.
\end{enumerate}

\section{Что остаётся программой}

Онтология, карта режимов и критерии отбора образуют дисциплинированную
программу, но не физическую теорию, пока отсутствует единый родительский закон
и независимое наблюдаемое следствие.

\section{Что считается частичным замыканием}

Отдельная формула, топологический класс или численно устойчивый коэффициент
могут быть закрыты локально. Их статус не переносится автоматически на всю
архитектуру.

\section{Условие перехода к следующему тому}

Следующий этап допустим, если первый том передаёт ему:
\begin{itemize}
 \item явный список определений;
 \item реестр гипотез и доказанных утверждений;
 \item критерии провала;
 \item перечень величин, запрещённых как скрытые входы;
 \item воспроизводимый порядок проверок.
\end{itemize}

Следующий крупный блок должен сформулировать почти-аксиоматический каркас,
леммы сохранения и квазидействие как условный генератор динамики.